% Auto-generated from 31-Deployment-Begins.md
% POV: Aisen | Landscape: Wasteland | Location: The Border Province

The first military convoy that moved Aisen toward Southeastern Territory occupied an entire day of travel through Covenant-controlled lands. The transformation from peaceful to militarized was gradual at first—military presence increasing incrementally in settlements they passed through—and then sudden as they approached the contested territories.

Checkpoints became more frequent. Military installations appeared with increasing regularity. Supply warehouses constructed for sustained occupation dominated the landscape. The infrastructure of military suppression was being built as they traveled through it.

Colonel Serath met Aisen's deployment group personally, which was unusually direct for someone of her rank to engage with junior officers. She greeted him with apparent casualness that masked careful evaluation—Aisen could see her tactical assessment framework hardened by decades of field operations.

"You're positioned as battalion intelligence liaison," she said as she showed him to the command post that had been established in a requisitioned government building in Ashfall's administrative district. "Your primary function is analyzing field intelligence from patrol operations and identifying patterns that help us understand provincial resistance structures."

"And my secondary function?" Aisen asked, testing.

Serath smiled slightly. "Your secondary function is communication liaison with higher command command structure. I need someone who can translate between field reality and command-level expectations. You appear to be that person based on your file notation."

The file notation she was referencing was Amara's insertion—the carefully coded personnel assessment that indicated Aisen was someone who could bridge operational perspectives.

"I'm also including you in command briefings," Serath continued. "Not all of them—some require classifications above your formal rank. But the ones where strategic pattern analysis is relevant, I want your perspective on what the data actually indicates versus what command narratives claim it indicates."

Aisen's first week in Southeastern Territory established operational patterns. He conducted analysis of patrol reports, identifying areas where resistance activity was concentrated. He coordinated with field intelligence specialists to understand what information was being gathered from interrogations and surveillance. He attended command briefings where strategic operations were being planned.

And every evening, he encoded messages into seemingly innocuous research documents that detailed:

- Military force distribution across the territory
- Character and tactical emphasis of military operations
- Identification of commanders who appeared sympathetic to constrained force approaches
- Supply chain vulnerabilities being exploited by resistance networks
- Projected military campaign timeline

The messages moved through covert channels established before deployment. Some went back to Kael's provincial network. Some went back to the Academy network through Vex. Some went to military officers who were themselves positioned as resistance sympathizers and needed current operational information to coordinate non-lethal suppression frameworks.

It was a remarkable balancing act—maintaining simultaneous allegiance to military command structure, provincial resistance networks, and military officers attempting to constrain the conflict to less brutal parameters.

The first major resistance operation occurred three weeks into the deployment. A military supply convoy was ambushed on a territory road. Forty soldiers were engaged by what appeared to be a coordinated resistance cell of approximately thirty fighters. The engagement lasted approximately ninety minutes before resistance fighters withdrew, having inflicted casualties and acquired military equipment.

The military casualty count was nine dead and fourteen wounded. The resistance coordination, based on field reports, appeared to have been sophisticated and well-executed.

Colonel Serath convened an immediate command meeting with section commanders, intelligence specialists, and tactical coordinators. Aisen was included in that meeting—a clear statement that she valued his analytical perspective on what the engagement meant.

"Pattern analysis," Serath said, directing her question to Aisen. "That level of coordination suggests what about the resistance organizational structure?"

"Regional-level command structure," Aisen said carefully, presenting his analysis while not revealing that he understood the resistance coordination because he was partially helping to facilitate it. "The ambush required knowledge of convoy timing, route, composition, and force concentration. That information level suggests coordination at command structure level, not individual cell. Provincial resistance movement has more organizational coherence than we initially assessed."

"Tactical implications?" another commander asked.

"They'll attempt more coordinated operations," Aisen said. "Probably targeting supply lines and isolated military positions. They're attempting to increase military operational cost without attempting direct confrontation. It's attrition strategy."

"How effective would that strategy be?" Serath asked.

"Moderately effective," Aisen responded. "If maintained long-term, it increases military supply requirements and personnel rotation pressures. Reduces operational efficiency. Not enough to defeat military force concentration but sufficient to make occupation logistically complex."

"Civilian casualty projections if we escalate response?" Serath asked pointedly, looking at Aisen directly.

This was the critical moment—other military commanders would have provided standard escalation frameworks. Aisen needed to provide honest analysis that supported constrained force approaches.

"If we maintain current military engagement parameters, I project civilian casualties will remain at current baseline rates—primarily as collateral from military operations and interrogation intensification," Aisen said. "If we escalate to area-level suppression tactics, projections increase significantly—twenty to thirty percent population casualty rate if we implement the full tactical framework command doctrine permits."

The room was silent. A thirty percent civilian casualty rate was... significant.

"Colonel," Serath said to the senior-most commander present, "do we have command authorization for population-level suppression tactics?"

"Theoretically yes," the senior commander replied. "But military doctrine notes that excessive civilian casualties can generate civilian population backlash that actually strengthens resistance effort long-term."

"Which means maintaining tactical restraint might actually support longer-term military objective achievement," Serath said. "Rather than undermining it."

"Correct assessment," Aisen agreed, which was understated—he was actually saying that mass civilian casualties would create conditions for stronger provincial resistance networks, which was entirely intentional from a resistance network perspective but strategically problematic from a military effectiveness perspective.

Colonel Serath's response to the ambush set the operational pattern for the following three months. Rather than escalating to maximum force suppression, she structured military operations around precise tactical engagement—targeted disruption of identified resistance activity without broad population impact.

It was exactly the kind of constrained force framework that Amara had advocated for. And it created conditions where the provincial resistance networks could maintain operational coherence without facing overwhelming military casualty-generation tactics.

But it also created tension with higher command structure, which wanted to see more aggressive military suppression and faster institutional replacement of provincial governance structures.

"Higher command is concerned that we're not conducting suppression operations with sufficient intensity," Serath told Aisen during a private briefing two months into the deployment. "They're pushing for timeline acceleration and casualty-rate increase."

"What are you saying?" Aisen asked carefully.

"I'm saying I need more detailed intelligence analysis about what specific military operations would maximize suppression effectiveness while minimizing civilian casualty rates," Serath said. "I need data that supports the position that tactical restraint is actually more effective than maximum force application."

It was a direct request for Aisen to use the intelligence position to help construct frameworks that supported Serath's operational philosophy. It was also, implicitly, a request for resistance network intelligence about what military operations would hurt provincial networks most effectively—creating justification for Serath to specifically avoid those operations.

"I can provide that analysis," Aisen said.

Over the following weeks, Aisen worked with Serath to construct detailed intelligence briefings that appeared to be purely military strategic analysis but were actually carefully constructed documents that:

1. Accurately identified the most significant threats to military occupation stability
2. Recommended precise tactical operations targeting those threats
3. Included enough plausible alternative analysis that Serath could claim she was exploring multiple engagement frameworks
4. Positioned civilian-minimization approaches as strategically sound rather than ethically motivated

The briefings were sophisticated enough that senior commanders accepted them as legitimate military analysis. They were also sophisticated enough that Serath could use them to justify operations that actually helped provincial resistance networks maintain organizational capacity.
